% ID: SEL-D0036-OB-E03
% ESTADO: cotejada. Segunda pasada visual y borrador compilado revisados.
% Fuente: Selectividad/matematicas_aplicadas/Ua1MATEMÁTICAS APLICADAS examen.pdf
% SHA-256: 77e0c295529ccd4ed502b5859d4011cb4b0e955a1e4795e570a9da20ae48293c
% Página: 2; opción B; ejercicio original 3.

El 30\% de los habitantes de una ciudad lee el diario A, el 13\% el diario B, y el 6\% ambos diarios.
\begin{enumerate}[label=\alph*)]
\item \textbf{(1.25 puntos)} \textquestiondown{}Qué porcentaje de habitantes de esta ciudad no lee ninguno de los diarios?
\item \textbf{(1.25 puntos)} Si se elige al azar un habitante de esta ciudad de entre los no lectores del diario B, \textquestiondown{}cuál es la probabilidad de que lea el diario A?
\end{enumerate}
